\chapter{User Interface Design}

\section{Two Dashboard Applications}

The project uses two dashboard applications because internal operators and
merchant users have different permissions, risks, and workflows.

\begin{table}[!htbp]
\centering
\caption{Dashboard split}
\begin{tabularx}{\textwidth}{L{0.24\textwidth}L{0.24\textwidth}Y}
\toprule
\textbf{UI} & \textbf{Audience} & \textbf{Purpose} \\
\midrule
Ops Dashboard & Internal Admin/Ops & Operate merchants, configure QR receiving accounts, investigate evidence, manage users, and recover failed workflows. \\
\bodyrowrule
Merchant Dashboard & Merchant dashboard users & Read merchant-scoped payment, refund, webhook, profile, credential, and analytics data. \\
\bottomrule
\end{tabularx}
\end{table}

This split follows the architectural access boundaries already established for
the system and keeps internal operations separate from merchant-facing
visibility. The result is a clearer user experience and stronger data
separation.

\section{Ops Dashboard UI}

The Ops Dashboard is designed for repeated operational work. It favors dense but
readable list and detail layouts, explicit lifecycle actions, semantic badges,
and direct access to operational evidence such as onboarding data, credentials,
VietQR receiving accounts, payments, refunds, webhook attempts, audit logs, and
reconciliation records.

The merchant detail page is also the control point for pilot QR readiness.
Internal users configure the merchant's receiving bank code, bank BIN, account
number, account name, template, and active or inactive state. This keeps QR
configuration in the same operational workflow as onboarding, credentials, and
activation.

\begin{figure}[!htbp]
\centering
\begin{subfigure}[t]{0.62\textwidth}
  \centering
  \safeimage[\linewidth]{ops-overview.png}{Ops Dashboard overview showing internal metrics and work queues.}
  \caption{Overview}
\end{subfigure}
\hfill
\begin{subfigure}[t]{0.32\textwidth}
  \centering
  \safeimage[\linewidth]{ops-merchant-detail.png}{Ops merchant detail with lifecycle status, credential metadata, portal users, onboarding evidence, recent payments, and audit.}
  \caption{Merchant detail}
\end{subfigure}
\caption{Ops Dashboard overview and merchant detail. The internal interface
supports both queue-level prioritization and case-level operational review
within one workflow.}
\end{figure}

\section{Merchant Dashboard UI}

The Merchant Dashboard is read-only except for local password change. It gives
merchant users an overview first, then deeper analytics and explorer pages for
support or business questions. Long values such as webhook URLs, QR payloads,
and JSON snippets are wrapped or scrollable to avoid layout overflow.

For payments that include generated VietQR data, the payment detail view can
show both the QR image and the underlying payload fields. Older or legacy
payments without a QR image still display their payload evidence in a full-width
layout so that the detail panel does not leave an empty visual gap.

\begin{figure}[!htbp]
\centering
\begin{subfigure}{0.48\textwidth}
  \centering
  \safeimage[\linewidth]{merchant-overview.png}{Merchant Overview with summary cards and compact charts.}
  \caption{Overview}
\end{subfigure}
\hfill
\begin{subfigure}{0.48\textwidth}
  \centering
  \safeimage[\linewidth]{merchant-analytics.png}{Merchant Analytics with interactive Recharts, tooltips, and data tables.}
  \caption{Analytics}
\end{subfigure}
\caption{Merchant Dashboard overview and analytics.}
\end{figure}

\begin{figure}[!htbp]
\centering
\begin{subfigure}{0.48\textwidth}
  \centering
  \safeimage[\linewidth]{merchant-payments.png}{Merchant payment explorer with filters, QR image, QR payload, and detail panel.}
  \caption{Payments}
\end{subfigure}
\hfill
\begin{subfigure}{0.48\textwidth}
  \centering
  \safeimage[\linewidth]{merchant-webhooks.png}{Merchant webhook explorer with event status and attempts.}
  \caption{Webhooks}
\end{subfigure}
\caption{Merchant explorer pages.}
\end{figure}

\section{Demo Merchant Checkout}

The demo merchant is a separate FastAPI application rather than another
gateway dashboard. Its setup panel accepts merchant integration credentials
once, while its server keeps the secret in memory and performs signed gateway
requests. The checkout displays the amount, transfer purpose, transaction ID,
countdown, and scannable VietQR image returned by the gateway.

The status timeline deliberately separates bank confirmation from merchant
notification. Pressing a simulation button sends a provider callback but leaves
the merchant order pending. The interface first states that the bank has
responded and that the merchant is waiting for the gateway. Only a valid signed
webhook received by the demo merchant changes the checkout to a terminal
success, failure, or expiration state.

\begin{figure}[!htbp]
\centering
\begin{subfigure}[t]{0.48\textwidth}
  \centering
  \safeimage[\linewidth]{demo-checkout-pending.png}{Demo merchant checkout displaying a pending VietQR payment, transfer details, timeline, and bank simulator controls.}
  \caption{Pending VietQR checkout}
\end{subfigure}
\hfill
\begin{subfigure}[t]{0.48\textwidth}
  \centering
  \safeimage[\linewidth]{demo-checkout-success.png}{Demo merchant checkout after the provider callback and signed gateway webhook, showing a successful payment result.}
  \caption{Webhook-confirmed success}
\end{subfigure}
\caption{The local demo merchant makes the asynchronous payment handoff visible.
The QR is generated by the gateway, while the terminal checkout result is based
only on the verified merchant webhook.}
\label{fig:demo-merchant-checkout}
\end{figure}

\section{Responsive And State Design}

Both dashboards treat asynchronous state as a first-class design concern:

\begin{itemize}
  \item loading, error, and empty states are explicit on data-fetching routes;
  \item long tables and side panels remain scrollable instead of forcing layout
  overflow;
  \item chart data is paired with readable labels and tables so exact values are
  not hidden behind hover-only interactions;
  \item session and status panels remain visually stable to reduce user
  confusion during multi-step workflows;
  \item the merchant login form uses explicit labels and accessible field
  associations.
  \item the demo checkout hides mutually exclusive empty and active states,
  keeps long transaction identifiers inside their columns, and collapses to a
  single-column layout without horizontal overflow on mobile screens.
\end{itemize}

Two smaller read-only pages complete the merchant-facing experience by exposing
merchant profile metadata and credential metadata without exposing raw secrets.

\begin{figure}[!htbp]
\centering
\begin{subfigure}{0.48\textwidth}
  \centering
  \safeimage[\linewidth]{merchant-profile.png}{Read-only merchant profile and integration metadata.}
  \caption{Profile}
\end{subfigure}
\hfill
\begin{subfigure}{0.48\textwidth}
  \centering
  \safeimage[\linewidth]{merchant-credentials.png}{Credential metadata without raw secret.}
  \caption{Credentials}
\end{subfigure}
\caption{Read-only merchant metadata pages.}
\end{figure}
